\chapter*{开始之前}

这本书讲的是一台机器：它重一千五百克，功耗二十瓦，和一只昏暗的灯泡差不多，却能认出人脸、学会语言、接住飞来的球。这就是你的大脑。我们要像工程师看一块陌生的电路板那样看它：上面有哪些元件，它们怎么连接，导线里跑的是什么信号，这套电路到底在计算什么。

你不需要公式，也不需要生物学知识。只需要好奇心，以及愿意想象导线、灯泡和开关。其余的，我们一路上边走边搭。

这本书还有另一个版本，即完整版，里面有方程、模型、习题和实验文献。两个版本的章节顺序相同，编号一一对应。如果哪一章让你着迷，就去读它的完整版：那里写明了每个论断从何而来，以及怎样自己去验证。在本书的网站上，每一章都可以用一个按钮在“简版”和“完整版”之间切换。

提醒一句。关于大脑，人类远没有全部弄清楚。我会坦白地说明，哪里是测量到的事实，哪里开始是猜测。这不是本书的缺点，反而是它最有意思的地方：许多章节以问题结尾，而这些问题目前还没有人知道答案。
